% GENERATED FILE. Edit canonical lessons, then run npm run generate:books.
% canonical-source-hash: fnv1a64:ca70a893f3153a37
% canonical-lessons: PA-C292-sillha, PA-C292-khushak, PA-C292-nami, PA-C292-sukkna, PA-C292-simmna

\chapter{Damp and Dry}
\label{ch:pa-a2-damp-and-dry}
\hlchaptermodality{292}

\begin{chapteropening}
\textbf{By the end of this chapter:} \emph{I can say whether the air and the walls are damp or dry.}

Complete the last lesson of chapter 292: I can say whether the air and the walls are damp or dry.
\end{chapteropening}

\section[\texorpdfstring{\pa{ਸਿੱਲ੍ਹਾ} (sillhā)}{ਸਿੱਲ੍ਹਾ (sillhā)}]{\pa{ਸਿੱਲ੍ਹਾ} (sillhā) — damp}

\label{lesson:PA-C292-sillha}



\begin{quote}
Before the new one: say the Punjabi for lukewarm, then the Punjabi for to cool something down.
\end{quote}

\subsection*{You'll want to know: \pa{ਸਿੱਲ੍ਹਾ}}
\textbf{\pa{ਸਿੱਲ੍ਹਾ}} — \emph{sillhā} — \textquotedblleft{}damp\textquotedblright{}.

\begin{quote}
\pa{ਬਰਸਾਤ} \pa{ਵਿੱਚ} \pa{ਕੰਧਾਂ} \pa{ਸਿੱਲ੍ਹੀਆਂ} \pa{ਹੁੰਦੀਆਂ} \pa{ਹਨ।}
\end{quote}

\begin{quote}
\emph{barsāt vich kandhāṁ sillhīāṁ hundīāṁ han.}
\end{quote}

\begin{quote}
In the rainy season the walls are damp.
\end{quote}

\begin{cousinweb}
\textbf{\pa{ਸਿੱਲ੍ਹਾ}} (\emph{sillhā}) is inherited from Sanskrit \emph{śītala}, \textquotedblleft{}cold\textquotedblright{}; CDIAL lists Punjabi \emph{sillhā}, \textquotedblleft{}damp\textquotedblright{}, under it (CDIAL 12487).
\end{cousinweb}

\begin{grammarlens}[title={What you've built}]
One more word for this chapter.
\end{grammarlens}

\subsection*{Your turn}
\begin{itemize}
\raggedright
  \item \emph{Say it:} \emph{sillhā}
  \item \emph{Say it:} \emph{sillhā}, once more
  \item \emph{Say it:} \emph{barsāt vich kandhāṁ sillhīāṁ hundīāṁ han.} — In the rainy season the walls are damp.
  \item \emph{Recall:} say the Punjabi for lukewarm, then the Punjabi for to cool something down, then say \emph{sillhā} again
\end{itemize}

\begin{tcolorbox}[breakable,colback=teal!4,colframe=teal!35!black,title={Before you move on}]
What does \pa{ਸਿੱਲ੍ਹਾ} mean? (\textquotedblleft{}Damp\textquotedblright{}.) Say it once more.
\end{tcolorbox}
\section[\texorpdfstring{\pa{ਖ਼ੁਸ਼ਕ} (khushak)}{ਖ਼ੁਸ਼ਕ (khushak)}]{\pa{ਖ਼ੁਸ਼ਕ} (khushak) — arid; dry in the air}

\label{lesson:PA-C292-khushak}



\begin{quote}
Before the new one: say the Punjabi for to cool something down, then the Punjabi for damp.
\end{quote}

\subsection*{You'll want to know: \pa{ਖ਼ੁਸ਼ਕ}}
\textbf{\pa{ਖ਼ੁਸ਼ਕ}} — \emph{khushak} — \textquotedblleft{}arid; dry in the air\textquotedblright{}.

\begin{quote}
\pa{ਸਿਆਲ} \pa{ਵਿੱਚ} \pa{ਹਵਾ} \pa{ਖ਼ੁਸ਼ਕ} \pa{ਹੁੰਦੀ} \pa{ਹੈ।}
\end{quote}

\begin{quote}
\emph{siāl vich havā khushak hundī hai.}
\end{quote}

\begin{quote}
In winter the air is dry.
\end{quote}

\begin{cousinweb}
\textbf{\pa{ਖ਼ੁਸ਼ਕ}} (\emph{khushak}) is a loan, usually traced to Persian \emph{khushk}, \textquotedblleft{}dry\textquotedblright{}.
\end{cousinweb}

\begin{grammarlens}[title={What you've built}]
One more word for this chapter.
\end{grammarlens}

\subsection*{Your turn}
\begin{itemize}
\raggedright
  \item \emph{Say it:} \emph{khushak}
  \item \emph{Say it:} \emph{khushak}, once more
  \item \emph{Say it:} \emph{siāl vich havā khushak hundī hai.} — In winter the air is dry.
  \item \emph{Recall:} say the Punjabi for to cool something down, then the Punjabi for damp, then say \emph{khushak} again
  \item \emph{Write it:} copy \textbf{\pa{ਖ਼ੁਸ਼ਕ}} once, checking the whole word
\end{itemize}

\begin{tcolorbox}[breakable,colback=teal!4,colframe=teal!35!black,title={Before you move on}]
What does \pa{ਖ਼ੁਸ਼ਕ} mean? (\textquotedblleft{}Arid; dry in the air\textquotedblright{}.) Say it once more.
\end{tcolorbox}
\section[\texorpdfstring{\pa{ਨਮੀ} (namī)}{ਨਮੀ (namī)}]{\pa{ਨਮੀ} (namī) — moisture; humidity}

\label{lesson:PA-C292-nami}



\begin{quote}
Before the new one: say the Punjabi for damp, then the Punjabi for arid.
\end{quote}

\subsection*{You'll want to know: \pa{ਨਮੀ}}
\textbf{\pa{ਨਮੀ}} — \emph{namī} — \textquotedblleft{}moisture; humidity\textquotedblright{}.

\begin{quote}
\pa{ਸਮੁੰਦਰ} \pa{ਦੇ} \pa{ਨੇੜੇ} \pa{ਹਵਾ} \pa{ਵਿੱਚ} \pa{ਬਹੁਤ} \pa{ਨਮੀ} \pa{ਹੁੰਦੀ} \pa{ਹੈ।}
\end{quote}

\begin{quote}
\emph{samundar de neṛe havā vich bahut namī hundī hai.}
\end{quote}

\begin{quote}
Near the sea there is a lot of moisture in the air.
\end{quote}

A feminine noun.

\begin{cousinweb}
\textbf{\pa{ਨਮੀ}} (\emph{namī}) is a loan, usually traced to Persian \emph{namī}, \textquotedblleft{}dampness\textquotedblright{}, built on \emph{nam}, \textquotedblleft{}damp\textquotedblright{}.
\end{cousinweb}

\begin{grammarlens}[title={What you've built}]
One more word for this chapter.
\end{grammarlens}

\subsection*{Your turn}
\begin{itemize}
\raggedright
  \item \emph{Say it:} \emph{namī}
  \item \emph{Say it:} \emph{namī}, once more
  \item \emph{Say it:} \emph{samundar de neṛe havā vich bahut namī hundī hai.} — Near the sea there is a lot of moisture in the air.
  \item \emph{Recall:} say the Punjabi for damp, then the Punjabi for arid, then say \emph{namī} again
  \item \emph{Write it:} copy \textbf{\pa{ਨਮੀ}} once, checking the whole word
\end{itemize}

\begin{tcolorbox}[breakable,colback=teal!4,colframe=teal!35!black,title={Before you move on}]
What does \pa{ਨਮੀ} mean? (\textquotedblleft{}Moisture; humidity\textquotedblright{}.) Say it once more.
\end{tcolorbox}
\section[\texorpdfstring{\pa{ਸੁੱਕਣਾ} (sukkṇā)}{ਸੁੱਕਣਾ (sukkṇā)}]{\pa{ਸੁੱਕਣਾ} (sukkṇā) — to dry up; to become dry}

\label{lesson:PA-C292-sukkna}



\begin{quote}
Before the new one: say the Punjabi for arid, then the Punjabi for moisture.
\end{quote}

\subsection*{You'll want to know: \pa{ਸੁੱਕਣਾ}}
\textbf{\pa{ਸੁੱਕਣਾ}} — \emph{sukkṇā} — \textquotedblleft{}to dry up; to become dry\textquotedblright{}.

\begin{quote}
\pa{ਧੁੱਪ} \pa{ਵਿੱਚ} \pa{ਕੱਪੜੇ} \pa{ਜਲਦੀ} \pa{ਸੁੱਕਦੇ} \pa{ਹਨ।}
\end{quote}

\begin{quote}
\emph{dhupp vich kappṛe jaldī sukkde han.}
\end{quote}

\begin{quote}
Clothes dry quickly in the sunshine.
\end{quote}

Past (it …): \pa{ਉਹ} \pa{ਕੱਪੜਾ} \pa{ਸੁੱਕਿਆ}, \emph{oh kappṛā sukkiā}.

\begin{cousinweb}
\textbf{\pa{ਸੁੱਕਣਾ}} (\emph{sukkṇā}) is inherited from a reconstructed \emph{*śuṣkati}, \textquotedblleft{}becomes dry\textquotedblright{}; CDIAL lists Punjabi \emph{sukkṇā}, \textquotedblleft{}dries up\textquotedblright{}, under it (CDIAL 12552).
\end{cousinweb}

\begin{grammarlens}[title={What you've built}]
One more word for this chapter.
\end{grammarlens}

\subsection*{Your turn}
\begin{itemize}
\raggedright
  \item \emph{Say it:} \emph{sukkṇā}
  \item \emph{Say it:} \emph{sukkṇā}, once more
  \item \emph{Say it:} \emph{dhupp vich kappṛe jaldī sukkde han.} — Clothes dry quickly in the sunshine.
  \item \emph{Recall:} say the Punjabi for arid, then the Punjabi for moisture, then say \emph{sukkṇā} again
  \item \emph{Write it:} copy \textbf{\pa{ਸੁੱਕਣਾ}} once, checking the whole word
\end{itemize}

\begin{tcolorbox}[breakable,colback=teal!4,colframe=teal!35!black,title={Before you move on}]
What does \pa{ਸੁੱਕਣਾ} mean? (\textquotedblleft{}To dry up; to become dry\textquotedblright{}.) Say it once more.
\end{tcolorbox}
\section[\texorpdfstring{\pa{ਸਿੰਮਣਾ} (simmṇā)}{ਸਿੰਮਣਾ (simmṇā)}]{\pa{ਸਿੰਮਣਾ} (simmṇā) — to seep; to ooze}

\label{lesson:PA-C292-simmna}



\begin{quote}
Before the new one: say the Punjabi for moisture, then the Punjabi for to dry up.
\end{quote}

\subsection*{You'll want to know: \pa{ਸਿੰਮਣਾ}}
\textbf{\pa{ਸਿੰਮਣਾ}} — \emph{simmṇā} — \textquotedblleft{}to seep; to ooze\textquotedblright{}.

\begin{quote}
\pa{ਮੀਂਹ} \pa{ਦਾ} \pa{ਪਾਣੀ} \pa{ਕੰਧ} \pa{ਵਿੱਚ} \pa{ਸਿੰਮਦਾ} \pa{ਹੈ।}
\end{quote}

\begin{quote}
\emph{mīṁh dā pāṇī kandh vich simmdā hai.}
\end{quote}

\begin{quote}
Rainwater seeps into the wall.
\end{quote}

Past (it …): \pa{ਉਹ} \pa{ਪਾਣੀ} \pa{ਸਿੰਮਿਆ}, \emph{oh pāṇī simmiā}.

\begin{cousinweb}
\textbf{\pa{ਸਿੰਮਣਾ}} (\emph{simmṇā}) is inherited from a reconstructed \emph{*simm}, \textquotedblleft{}ooze\textquotedblright{}; CDIAL lists Punjabi \emph{simmṇā}, \textquotedblleft{}to ooze\textquotedblright{}, under it (CDIAL 13419).
\end{cousinweb}

\begin{grammarlens}[title={What you've built}]
One more word for this chapter.
\end{grammarlens}

\subsection*{Your turn}
\begin{itemize}
\raggedright
  \item \emph{Say it:} \emph{simmṇā}
  \item \emph{Say it:} \emph{simmṇā}, once more
  \item \emph{Say it:} \emph{mīṁh dā pāṇī kandh vich simmdā hai.} — Rainwater seeps into the wall.
  \item \emph{Recall:} say the Punjabi for moisture, then the Punjabi for to dry up, then say \emph{simmṇā} again
  \item \emph{Write it:} copy \textbf{\pa{ਸਿੰਮਣਾ}} once, checking the whole word
\end{itemize}

\begin{tcolorbox}[breakable,colback=teal!4,colframe=teal!35!black,title={Before you move on}]
What does \pa{ਸਿੰਮਣਾ} mean? (\textquotedblleft{}To seep; to ooze\textquotedblright{}.) Say it once more.
\end{tcolorbox}
